\section{Two-key rectangles and cross-key holonomy}
\label{sec:cross-key}

The shortest mixed closed word is \(ER\).  By
\cref{cor:first-mixed-cross-key}, it necessarily uses two distinct complete
native keys.  We now write its four corners explicitly.

\subsection{Exact rectangle formula}

For keys \(i,j\) and rows \(m,n\), define
\begin{equation}
 \mathcal H_{i,j}(m,n)
 :=
 \overline{f_i(m)}f_i(n)
 \overline{f_j(n)}f_j(m).
 \label{eq:rectangle-holonomy}
\end{equation}
It vanishes when one of the four key-row corners is absent.
Throughout this section, a residual-relation indicator is also defined to
be zero if either amplitude named by that indicator vanishes.

\begin{theorem}[Cross-key rectangle identity]
\label{thm:cross-key-rectangle}
For \(i\ne j\),
\begin{equation}
\boxed{
\begin{aligned}
\tr(E_iR_j)
={}&
2\operatorname{Re}
\sum_{m<n}
\one_{\{s_i(m)\ne s_i(n)\}}
\one_{\{s_j(m)=s_j(n)\}}
\mathcal H_{i,j}(m,n).
\end{aligned}}
\label{eq:cross-key-rectangle}
\end{equation}
Consequently,
\begin{equation}
\boxed{
\tr(ER)
=
2\operatorname{Re}
\sum_{\substack{i\ne j\\m<n}}
\one_{\{s_i(m)\ne s_i(n)\}}
\one_{\{s_j(m)=s_j(n)\}}
\mathcal H_{i,j}(m,n).}
\label{eq:full-cross-key-rectangle}
\end{equation}
\end{theorem}

\begin{proof}
By definition,
\[
 \tr(E_iR_j)
 =
 \sum_{m,n}(E_i)_{mn}(R_j)_{nm}.
\]
The diagonal terms vanish.  For \(m\ne n\), substitute
\[
 (E_i)_{mn}
 =
 \one_{\{s_i(m)\ne s_i(n)\}}
 \overline{f_i(m)}f_i(n)
\]
and
\[
 (R_j)_{nm}
 =
 \one_{\{s_j(n)=s_j(m)\}}
 \overline{f_j(n)}f_j(m).
\]
The ordered terms \((m,n)\) and \((n,m)\) are complex conjugates, giving
\eqref{eq:cross-key-rectangle}.  Sum over \(i\ne j\) and use
\eqref{eq:first-mixed-cross-key}.
\end{proof}

The key \(i\) carries the unequal-residual edge and \(j\) carries the
same-residual edge.  The roles are not interchangeable in
\eqref{eq:full-cross-key-rectangle}; the term with \((j,i)\) is present
only if the residual relations are reversed at those keys.

\subsection{Gauge invariance}

\begin{proposition}[Rectangle holonomy is gauge invariant]
\label{prop:holonomy-gauge-invariant}
Let \(u_m,v_i\in\mathbb C\) have modulus one and replace
\[
 f_i(m)\longmapsto u_mv_if_i(m).
\]
Then every \(\mathcal H_{i,j}(m,n)\) is unchanged.
\end{proposition}

\begin{proof}
In \eqref{eq:rectangle-holonomy}, the row phases occur as
\[
 \overline{u_m}u_n\overline{u_n}u_m=1
\]
and the key phases as
\[
 \overline{v_i}v_i\overline{v_j}v_j=1.
\]
\end{proof}

Thus the rectangle phase cannot be removed by moving a unit phase between
source rows and outputs.  It is a genuine four-corner quantity.  Support
and coefficient magnitudes do not determine it.

\subsection{The residual sign on a rectangle}

Let \(\beta_{i,m}\) be the literal coefficient of the unique atom at
\((i,m)\), extended by zero.  The row-gauged amplitude is
\[
 f_i(m)
 =
 \mu(n_{i,m})\frac{\beta_{i,m}}{\sqrt{W_m}},
\qquad
 n_{i,m}=m\,j(i)+h_0.
\]

\begin{proposition}[Four target signs reduce to one erasure pair]
\label{prop:rectangle-sign-reduction}
On a rectangle contributing to
\eqref{eq:cross-key-rectangle},
\begin{equation}
\boxed{
\mathcal H_{i,j}(m,n)
=
\mu(s_i(m))\mu(s_i(n))
\frac{
\overline{\beta_{i,m}}\beta_{i,n}
\overline{\beta_{j,n}}\beta_{j,m}
}{W_mW_n}.}
\label{eq:rectangle-sign-reduction}
\end{equation}
\end{proposition}

\begin{proof}
The product of the four target signs is
\[
 \mu(n_{i,m})\mu(n_{i,n})
 \mu(n_{j,n})\mu(n_{j,m}).
\]
Use
\[
 \mu(n_{k,m})=-\mu(d_m)\mu(s_k(m)).
\]
The four minus signs cancel, and each source-row sign occurs twice.  At the
same-residual key \(j\),
\(\mu(s_j(m))\mu(s_j(n))=1\).  The two residual signs at the erasure key
\(i\) remain, proving the formula.
\end{proof}

Equation \eqref{eq:rectangle-sign-reduction} is the minimal instance of the
residual-boundary theorem.  It also shows why a one-key estimate cannot
see the first mixed response: the sign-bearing unequal-residual edge and
its same-residual return occupy different complete native keys.

\subsection{The direct-cell two-time rectangle}

In the identical-support direct-cell specialization, let \(c\) collect a
fixed side and external cell, and let
\[
 \mathfrak m_m^c(t)
\]
be the actual missing-time indicator.  Put
\[
 \chi_E^{mn}(t)=\one_{\{s_m(t)\ne s_n(t)\}},
 \qquad
 \chi_R^{mn}(t)=\one_{\{s_m(t)=s_n(t)\}}.
\]
After summing the two independent orthogonal \(\gamma\)-coordinates
exactly, \eqref{eq:full-cross-key-rectangle} becomes
\begin{equation}
\boxed{
\begin{aligned}
\tr(ER)
={}&
2\operatorname{Re}
\sum_{m<n}\frac1{W_mW_n}
\sum_{\substack{c,t\\c',u}}
\mathfrak m_m^c(t)\mathfrak m_n^c(t)
\mathfrak m_n^{c'}(u)\mathfrak m_m^{c'}(u)
\\
&\times
\chi_E^{mn}(t)\chi_R^{mn}(u)
\mu(s_m(t))\mu(s_n(t))
\\
&\times
\overline{a_m^c(t)}a_n^c(t)
\overline{a_n^{c'}(u)}a_m^{c'}(u)
\Gamma_{mn}^c(t)\Gamma_{nm}^{c'}(u).
\end{aligned}}
\label{eq:direct-cell-two-key-rectangle}
\end{equation}
The two times \(t\) and \(u\), two sides/cells \(c,c'\), and the two
\(\gamma\)-sums are independent.  Formula
\eqref{eq:direct-cell-two-key-rectangle} is invalid outside the declared
specialization unless the full \(\gamma\)-dependent first-failure
indicators are restored.

\subsection{Two determinant edges}

For the \(i\)-edge at time \(t=j(i)\), write
\[
 p_{i,m}r_{i,m}=mt+h_0,\qquad
 p_{i,n}r_{i,n}=nt+h_0.
\]
Then
\begin{equation}
 n p_{i,m}r_{i,m}-m p_{i,n}r_{i,n}
 =h_0(n-m).
 \label{eq:i-edge-determinant}
\end{equation}
The same identity holds at key \(j\), generally with another time.  The
rectangle conditions are
\begin{align}
 d_mr_{i,m}&\ne d_nr_{i,n},
 \label{eq:E-determinant-condition}\\
 d_mr_{j,m}&=d_nr_{j,n}.
 \label{eq:R-determinant-condition}
\end{align}
Thus the first mixed trace is a pair of determinant solutions sharing two
source rows, with unequal residual labels on one horizontal edge and equal
residual labels on the other.  The determinant identities are exact
reindexings; they do not pair the signs in
\eqref{eq:rectangle-sign-reduction}.
